\begin{figure}[!htbp]
\centering
\begin{tikzpicture}[
  x=1cm,y=1cm,
  every node/.style={font=\small},
  common/.style={draw=black!55,fill=black!16,line width=.5pt},
  marked/.style={draw=BurntOrange,fill=BurntOrange!30,line width=.9pt}
]
  \node[anchor=east] at (0.8,1.35) {$T_0$};
  \node[anchor=east] at (0.8,-0.35) {$T_1$};

  \foreach \y in {1.35,-0.35} {
    \draw[common] (1.55,\y-.16) rectangle (1.87,\y+.16);
    \draw[common] (2.17,\y-.16) rectangle (2.49,\y+.16);
    \node at (2.94,\y) {$\cdots$};
    \draw[common] (3.40,\y-.16) rectangle (3.72,\y+.16);
    \node at (2.64,\y-.51) {$T_*$};
  }

  \draw[marked] (4.84,1.19) rectangle (5.16,1.51);
  \node at (5,0.84) {$t(0)$};

  \draw[draw=black!45,dashed,line width=.6pt]
    (4.84,-.51) rectangle (5.16,-.19);
  \draw[-{Latex[length=2.2mm]},line width=.8pt,BurntOrange]
    (5.40,-.35) -- node[above=3pt,text=black] {$mw$} (9.36,-.35);
  \draw[marked] (9.60,-.51) rectangle (9.92,-.19);
  \node at (9.76,-.86) {$t(0)+mw$};
\end{tikzpicture}
\caption{Schematic of the marked-point switch, without depicting the
actual positions or scale.  The gray marks represent the same common
set $T_*$ in both rows.  The marked point moves from $t(0)$ to
$t(0)+mw$; the dashed outline shows its former location.  The marks
may also denote unit-cube thickenings.  In a grid assembly this switch
gives the kernel-invariance identity \eqref{geom:kernel-invariance}.}
\label{geom:fig-switch}
\end{figure}
